\begin{proof}
Let \(a,b\ge 0\) and let \(p,q>1\) satisfy \(\dfrac1p+\dfrac1q=1\).

\medskip
\noindent\textbf{Degenerate cases.}
If \(a=0\) or \(b=0\) then \(ab=0\) while the right–hand side  
\(\dfrac{a^{p}}{p}+\dfrac{b^{q}}{q}\) is non‑negative, so the inequality holds.
Hence we may assume
\[
a>0,\qquad b>0 .
\]

\medskip
\noindent\textbf{A convex function.}
Define \(\displaystyle \phi(x)=\frac{x^{p}}{p}\) for \(x\ge 0\).
Then
\[
\phi'(x)=x^{p-1},\qquad 
\phi''(x)=(p-1)x^{p-2}>0\quad (p>1,\;x>0),
\]
so \(\phi\) is convex on \([0,\infty)\).  
For a convex differentiable function the supporting‑line inequality holds:
\begin{equation}
\phi(a)\ge \phi(x_{0})+\phi'(x_{0})(a-x_{0})
\qquad(\forall a\ge0,\;x_{0}>0). \tag{1}
\end{equation}

\medskip
\noindent\textbf{Choice of the tangency point.}
Take \(x_{0}=b^{\,q-1}\).  
Since the exponents are conjugate,
\[
q=\frac{p}{p-1},\qquad q-1=\frac{1}{p-1},
\]
and therefore \((q-1)(p-1)=1\). Consequently
\begin{align}
\phi'(x_{0})&=(b^{\,q-1})^{p-1}=b^{(q-1)(p-1)}=b, \tag{2}\\[4pt]
\phi(x_{0})&=\frac{(b^{\,q-1})^{p}}{p}
          =\frac{b^{p(q-1)}}{p}
          =\frac{b^{q}}{p}, \tag{3}
\end{align}
because \(p(q-1)=q\).

\medskip
\noindent\textbf{Applying (1).}
Insert \(x_{0}=b^{\,q-1}\) together with (2) and (3) into (1):
\[
\frac{a^{p}}{p}
   \ge \frac{b^{q}}{p}+b\bigl(a-b^{\,q-1}\bigr).
\]
Expanding the right–hand side yields
\[
\frac{a^{p}}{p}\ge \frac{b^{q}}{p}+ab-b^{q}.
\]
Rearranging,
\[
\frac{a^{p}}{p}+\Bigl(1-\frac1p\Bigr)b^{q}\ge ab.
\]
Since \(1-\dfrac1p=\dfrac{p-1}{p}=\dfrac1q\) (by \(\frac1p+\frac1q=1\)), we obtain
\begin{equation}
\frac{a^{p}}{p}+\frac{b^{q}}{q}\ge ab. \tag{4}
\end{equation}
Thus Young’s inequality holds for all \(a,b>0\).

\medskip
\noindent\textbf{Equality condition.}
In (1) equality occurs iff \(a=x_{0}\). With our choice of \(x_{0}\) this means
\(a=b^{\,q-1}\). Raising both sides to the power \(p\) and using \(p(q-1)=q\) gives
\(a^{p}=b^{q}\). Conversely, if \(a^{p}=b^{q}\) then \(a=b^{\,q-1}\) and each step
above is an equality. The trivial case \(a=b=0\) also gives equality. Hence
\[
ab=\frac{a^{p}}{p}+\frac{b^{q}}{q}
\quad\Longleftrightarrow\quad a^{p}=b^{q}.
\]

\medskip
\noindent\textbf{Conclusion.}
Combining the zero case with the argument for \(a,b>0\) we have proved Young’s
inequality
\[
ab\le\frac{a^{p}}{p}+\frac{b^{q}}{q},
\]
with equality precisely when \(a^{p}=b^{q}\) (including \(a=b=0\)). \(\square\)
\end{proof}